\section{Related Work and Positioning}
Anomaly-based data prioritization has spacecraft precedent. Finley and
colleagues investigate inexpensive anomaly detection for prioritizing in situ
observations~\citep{finley2023prioritization,finley2024generalized}. Our focus is engineering-fault context under explicit recording budgets, rather than
claiming anomaly-guided downlink selection itself as a contribution. The evaluation therefore concerns retained evidence under a declared byte
contract, without claiming anomaly-guided downlink selection as a new contribution.

Joint compression and anomaly detection has also been studied in networked sensor
systems~\citep{elsayed2023compression}. Our first implementation uses transparent
raw and extrema records rather than a learned codec. This makes metadata costs
and retained information auditable, while leaving stronger compression baselines
for the full evaluation.

ESA-ADB supplies annotated mission telemetry and a hierarchical evaluation
framework~\citep{kotowski2024esa}. We use its native Mission1 observations for retention experiments. Detector performance on a complete archive and evidence
preservation after budgeted capture are different questions. The retention-specific outcomes do not replace its official hierarchical
anomaly-detection metrics, which are not reported here.
Reservoir sampling provides a classical detector-independent streaming control
when the final observation count is unknown~\citep{vitter1985reservoir}. We use
Algorithm R rather than Algorithm Z and add randomized temporal phase to test
alignment sensitivity. These are established controls, not claimed innovations.
The present comparisons do not establish novelty or superiority relative to
stronger adaptive acquisition, compression, or established onboard fault recorders.
